\chapter{Data Visualization in R}
\label{chap:visualization}

\chapterguide
  {You should be able to load and inspect data frames.}
  {create and interpret pie charts, box plots, histograms, scatterplots, bar charts, and line graphs.}

\section{Using built-in datasets}

List built-in datasets, then load and inspect \texttt{mtcars}:

\begin{lstlisting}[style=dsr]
data()

data(mtcars)
print(head(mtcars, 10))

input <- mtcars[, c("mpg", "cyl")]
print(head(input))
\end{lstlisting}

The workflow is: load, inspect, select, then plot.

\section{Built-in graphics demonstrations}

\begin{lstlisting}[style=dsr]
demo(graphics)
demo(image)
\end{lstlisting}

Press Enter to advance through each demonstration.

\section{Pie charts}

A basic pie chart uses values and labels:

\begin{lstlisting}[style=dsr]
x <- c(21, 62, 10, 53)
labels <- c("London", "New York", "Singapore", "Mumbai")

pie(x, labels)
\end{lstlisting}

Save the plot by opening a graphics device and closing it afterward:

\begin{lstlisting}[style=dsr]
png(file = "city.png")
pie(x, labels)
dev.off()
\end{lstlisting}

\subsection{Title, colors, percentages, and legend}

\begin{lstlisting}[style=dsr]
pie(
  x,
  labels,
  main = "City pie chart",
  col = rainbow(length(x))
)

piepercent <- round(100 * x / sum(x), 1)

pie(
  x,
  labels = piepercent,
  main = "City pie chart",
  col = rainbow(length(x))
)

legend(
  "topright",
  c("London", "New York", "Singapore", "Mumbai"),
  cex = 0.8,
  fill = rainbow(length(x))
)
\end{lstlisting}

The labels show each value as a percentage of the total.

\begin{utpfig}[Pie chart]{The chart the code above produces. Slice angle encodes
each city's share of the total 146.}{fig:pie}
% Generated by tools/make-figures.py -- do not edit by hand.
% Regenerate with: python3 tools/make-figures.py
\begin{tikzpicture}[scale=1.0]
\fill[UTPBlue!75,draw=white,line width=1pt] (0,0) -- (90.00:2.1) arc[start angle=90.00, end angle=38.22, radius=2.1] -- cycle;
\node[font=\scriptsize\bfseries\color{white}] at (64.11:1.35) {14.4\%};
\node[font=\scriptsize\color{UTPInk}] at (64.11:2.62) {London};
\fill[UTPTeal!75,draw=white,line width=1pt] (0,0) -- (38.22:2.1) arc[start angle=38.22, end angle=-114.66, radius=2.1] -- cycle;
\node[font=\scriptsize\bfseries\color{white}] at (-38.22:1.35) {42.5\%};
\node[font=\scriptsize\color{UTPInk}] at (-38.22:2.62) {New York};
\fill[UTPOrange!75,draw=white,line width=1pt] (0,0) -- (-114.66:2.1) arc[start angle=-114.66, end angle=-139.32, radius=2.1] -- cycle;
\node[font=\scriptsize\bfseries\color{white}] at (-126.99:1.35) {6.8\%};
\node[font=\scriptsize\color{UTPInk}] at (-126.99:2.62) {Singapore};
\fill[UTPRed!75,draw=white,line width=1pt] (0,0) -- (-139.32:2.1) arc[start angle=-139.32, end angle=-270.00, radius=2.1] -- cycle;
\node[font=\scriptsize\bfseries\color{white}] at (-204.66:1.35) {36.3\%};
\node[font=\scriptsize\color{UTPInk}] at (-204.66:2.62) {Mumbai};
\end{tikzpicture}

\end{utpfig}

\begin{pitfall}
Comparing Mumbai (36.3\,\%) with New York (42.5\,\%) means comparing two
angles, which people do poorly. A bar chart of the same four numbers would
be read correctly at a glance. Use a pie only when the point is
\emph{share of a whole} and the categories number about five or fewer.
\end{pitfall}

\subsection{3D pie chart}

The optional \texttt{plotrix} example requires defined labels:

\begin{lstlisting}[style=dsr]
library(plotrix)

pie3D(
  x,
  labels = labels,
  explode = 0.1,
  main = "Pie Chart of Countries"
)
\end{lstlisting}

\begin{pitfall}
The source uses an undefined object, \texttt{lbl}; the working example reuses \texttt{labels}. Prefer a clear 2D chart when depth adds no information.
\end{pitfall}

\section{Box plots}

Compare mileage across cylinder groups with formula notation:

\begin{lstlisting}[style=dsr]
boxplot(
  mpg ~ cyl,
  data = mtcars,
  xlab = "Number of Cylinders",
  ylab = "Miles Per Gallon",
  main = "Mileage Data"
)
\end{lstlisting}

A more customized version adds notches, variable widths, colors, and supplied group names:

\begin{lstlisting}[style=dsr]
boxplot(
  mpg ~ cyl,
  data = mtcars,
  xlab = "Number of Cylinders",
  ylab = "Miles Per Gallon",
  main = "Mileage Data",
  notch = TRUE,
  varwidth = TRUE,
  col = c("green", "yellow", "purple"),
  names = c("4 cylinders", "6 cylinders", "8 cylinders")
)
\end{lstlisting}

\begin{utpfig}[Box plot]{\texttt{mpg} by cylinder count, from the 32 rows of
\texttt{mtcars}. The box spans the interquartile range, the line inside it is
the median, and the whiskers reach the furthest point within 1.5 IQR. Points
beyond that are drawn individually.}{fig:boxplot}
% Generated by tools/make-figures.py -- do not edit by hand.
% Regenerate with: python3 tools/make-figures.py
\begin{tikzpicture}
\begin{axis}[
width=0.62\linewidth, height=5.4cm, tick align=outside, tick pos=left,
  grid=major, grid style={draw=UTPMuted!20},
  axis line style={draw=UTPMuted!70},
  boxplot/draw direction=y,
  xtick={1,2,3}, xticklabels={4 cylinders,6 cylinders,8 cylinders},
  xlabel={Number of Cylinders}, ylabel={Miles Per Gallon},
  ymin=8, ymax=36, xticklabel style={font=\scriptsize},
]
\addplot+[boxplot prepared={lower whisker=21.40, lower quartile=22.800, median=26.00, upper quartile=30.400, upper whisker=33.90, draw position=1}, color=UTPNavy, fill=UTPGreen!35, solid, mark options={color=UTPRed}] coordinates {};
\addplot+[boxplot prepared={lower whisker=17.80, lower quartile=18.650, median=19.70, upper quartile=21.000, upper whisker=21.40, draw position=2}, color=UTPNavy, fill=UTPGold!35, solid, mark options={color=UTPRed}] coordinates {};
\addplot+[boxplot prepared={lower whisker=13.30, lower quartile=14.400, median=15.20, upper quartile=16.250, upper whisker=18.70, draw position=3}, color=UTPNavy, fill=UTPPurple!35, solid, mark options={color=UTPRed}] coordinates {(3,10.4) (3,10.4) (3,19.2)};
\end{axis}
\end{tikzpicture}

\end{utpfig}

\begin{keyidea}
The formula \texttt{mpg \textasciitilde{} cyl} reads ``mpg \emph{by} cyl''.
It splits one numeric column into groups defined by another column, which is
why a single \texttt{boxplot()} call produces three boxes rather than one.
\end{keyidea}

\section{Histograms}

\begin{lstlisting}[style=dsr]
v <- c(9, 13, 21, 8, 36, 22, 12, 41, 31, 33, 19)

hist(
  v,
  xlab = "Weight",
  col = "yellow",
  border = "blue"
)
\end{lstlisting}

Add axis limits and a requested break count:

\begin{lstlisting}[style=dsr]
hist(
  v,
  xlab = "Weight",
  col = "green",
  border = "red",
  xlim = c(0, 50),
  ylim = c(0, 5),
  breaks = 5
)
\end{lstlisting}

\begin{utpfig}[Histogram]{The 11 values binned into five intervals of width 10.
Bar height is a count, so the bars touch: the horizontal axis is continuous,
unlike a bar chart's.}{fig:histogram}
% Generated by tools/make-figures.py -- do not edit by hand.
% Regenerate with: python3 tools/make-figures.py
\begin{tikzpicture}
\begin{axis}[
width=0.62\linewidth, height=5.0cm, tick align=outside, tick pos=left,
  grid=major, grid style={draw=UTPMuted!20},
  axis line style={draw=UTPMuted!70},
  ybar interval=1,
  xlabel={Weight}, ylabel={Frequency},
  xmin=0, xmax=50, ymin=0, ymax=5,
  xtick={0,10,20,30,40,50}, ytick={0,1,2,3,4,5},
]
\addplot+[fill=UTPGreen!45, draw=UTPRed, solid, mark=none]
  coordinates {(0,2) (10,3) (20,2) (30,3) (40,1) (50,0)};
\end{axis}
\end{tikzpicture}

\end{utpfig}

\begin{pitfall}
The source prints the invalid argument \texttt{n breaks = 5}; the correct argument is \texttt{breaks = 5}.

\texttt{breaks} is a suggestion, not an instruction. R applies
\texttt{pretty()} to it and may return a different number of bins than you
asked for. Check the result rather than assuming.
\end{pitfall}

\section{Scatterplots}

Lab 8a selects weight and mileage from \texttt{mtcars}:

\begin{lstlisting}[style=dsr]
input <- mtcars[, c("wt", "mpg")]

plot(
  x = input$wt,
  y = input$mpg,
  xlab = "Weight",
  ylab = "Mileage",
  xlim = c(2.5, 5),
  ylim = c(15, 30),
  main = "Weight versus mileage"
)
\end{lstlisting}

\begin{utpfig}[Scatterplot]{Weight against mileage, drawn with the axis limits
the lab specifies. The downward drift is the finding: heavier cars return fewer
miles per gallon.}{fig:scatter}
% Generated by tools/make-figures.py -- do not edit by hand.
% Regenerate with: python3 tools/make-figures.py
\begin{tikzpicture}
\begin{axis}[
width=0.62\linewidth, height=5.0cm, tick align=outside, tick pos=left,
  grid=major, grid style={draw=UTPMuted!20},
  axis line style={draw=UTPMuted!70},
  xlabel={Weight}, ylabel={Mileage},
  xmin=2.5, xmax=5, ymin=15, ymax=30,
]
\addplot+[only marks, mark=o, mark size=1.9pt, draw=UTPBlue, solid]
  coordinates {(2.620,21.0) (2.875,21.0) (3.215,21.4) (3.440,18.7) (3.460,18.1) (3.190,24.4) (3.150,22.8) (3.440,19.2) (3.440,17.8) (4.070,16.4) (3.730,17.3) (3.780,15.2) (3.520,15.5) (3.435,15.2) (3.845,19.2) (3.170,15.8) (2.770,19.7) (3.570,15.0) (2.780,21.4)};
\end{axis}
\end{tikzpicture}
% 13 of the 32 mtcars rows fall outside the lab's xlim/ylim and are
% therefore not drawn -- that is the point of the accompanying pitfall.

\end{utpfig}

\begin{pitfall}
\texttt{xlim} and \texttt{ylim} do not rescale the plot --- they discard
anything outside the window, silently. These limits drop 13 of the 32
\texttt{mtcars} rows, including every car under 2.5 tons and the three
returning over 30\,mpg, which are exactly the points that make the trend
clearest. Plot without limits first, then narrow them only if you can say
what you are excluding.
\end{pitfall}

A scatterplot matrix expands the idea to four variables:

\begin{lstlisting}[style=dsr]
pairs(
  ~ wt + mpg + disp + cyl,
  data = mtcars,
  main = "Scatterplot Matrix"
)
\end{lstlisting}

\section{Bar charts}

A basic bar chart is:

\begin{lstlisting}[style=dsr]
H <- c(7, 12, 28, 3, 41)
barplot(H)
\end{lstlisting}

The customized version supplies month names, labels, title, fill color, and border color:

\begin{lstlisting}[style=dsr]
H <- c(7, 12, 28, 3, 41)
M <- c("Mar", "Apr", "May", "Jun", "Jul")

barplot(
  H,
  names.arg = M,
  xlab = "Month",
  ylab = "Revenue",
  col = "blue",
  main = "Revenue chart",
  border = "red"
)
\end{lstlisting}

\begin{utpfig}[Bar chart]{The same five numbers as \figref{fig:linegraph}, drawn
as bars. Bars are separated because the horizontal axis is categorical: the
months have no numeric distance between them.}{fig:barchart}
% Generated by tools/make-figures.py -- do not edit by hand.
% Regenerate with: python3 tools/make-figures.py
\begin{tikzpicture}
\begin{axis}[
width=0.60\linewidth, height=4.6cm, tick align=outside, tick pos=left,
  grid=major, grid style={draw=UTPMuted!20},
  axis line style={draw=UTPMuted!70},
  ybar, bar width=15pt,
  symbolic x coords={Mar,Apr,May,Jun,Jul}, xtick=data,
  xlabel={Month}, ylabel={Revenue}, ymin=0, ymax=45,
]
\addplot+[fill=UTPBlue!70, draw=UTPRed, solid] coordinates {(Mar,7) (Apr,12) (May,28) (Jun,3) (Jul,41)};
\end{axis}
\end{tikzpicture}

\end{utpfig}

\subsection{Stacked matrix input}

Lab 8a constructs a 3-by-5 matrix of values and passes it to \texttt{barplot()} with a region legend:

\begin{lstlisting}[style=dsr]
colors <- c("green", "orange", "brown")
months <- c("Mar", "Apr", "May", "Jun", "Jul")
regions <- c("East", "West", "North")

Values <- matrix(
  c(2,9,3,11,9, 4,8,7,3,12, 5,2,8,10,11),
  nrow = 3,
  ncol = 5,
  byrow = TRUE
)

barplot(
  Values,
  main = "Total revenue",
  names.arg = months,
  xlab = "Month",
  ylab = "Revenue",
  col = colors
)

legend("topleft", regions, cex = 1.3, fill = colors)
\end{lstlisting}

\begin{utpfig}[Stacked bar chart]{Three regions stacked per month. Total height
is easy to compare; the middle segment is not, because it starts at a different
place in every column.}{fig:stacked}
% Generated by tools/make-figures.py -- do not edit by hand.
% Regenerate with: python3 tools/make-figures.py
\begin{tikzpicture}
\begin{axis}[
width=0.60\linewidth, height=4.8cm, tick align=outside, tick pos=left,
  grid=major, grid style={draw=UTPMuted!20},
  axis line style={draw=UTPMuted!70},
  ybar stacked, bar width=15pt,
  symbolic x coords={Mar,Apr,May,Jun,Jul}, xtick=data,
  xlabel={Month}, ylabel={Revenue}, ymin=0,
  legend style={at={(0.02,0.98)}, anchor=north west, font=\scriptsize},
  legend cell align=left,
]
\addplot+[ybar stacked, fill=UTPGreen!60, draw=UTPNavy, solid] coordinates {(Mar,2) (Apr,9) (May,3) (Jun,11) (Jul,9)};
\addplot+[ybar stacked, fill=UTPOrange!60, draw=UTPNavy, solid] coordinates {(Mar,4) (Apr,8) (May,7) (Jun,3) (Jul,12)};
\addplot+[ybar stacked, fill=UTPRed!60, draw=UTPNavy, solid] coordinates {(Mar,5) (Apr,2) (May,8) (Jun,10) (Jul,11)};
\legend{East,West,North}
\end{axis}
\end{tikzpicture}

\end{utpfig}

\begin{pitfall}
Only the bottom segment of a stacked bar shares a baseline, so only it can be
compared across columns by eye. If the question is about West rather than the
total, use grouped bars (\texttt{beside = TRUE}) instead.
\end{pitfall}

\section{Line graphs}

\begin{lstlisting}[style=dsr]
v <- c(7, 12, 28, 3, 41)
plot(v, type = "o")
\end{lstlisting}

Add a title, labels, and color:

\begin{lstlisting}[style=dsr]
plot(
  v,
  type = "o",
  col = "red",
  xlab = "Month",
  ylab = "Rainfall",
  main = "Rainfall chart"
)
\end{lstlisting}

A second series is added using \texttt{lines()}:

\begin{lstlisting}[style=dsr]
t <- c(14, 7, 6, 19, 3)

plot(
  v,
  type = "o",
  col = "red",
  xlab = "Month",
  ylab = "Rainfall",
  main = "Rainfall chart"
)
lines(t, type = "o", col = "blue")
\end{lstlisting}

\begin{utpfig}[Line graph]{Two series on one pair of axes. \texttt{plot()} draws
the first and fixes the axis limits; \texttt{lines()} adds to what is already
there.}{fig:linegraph}
% Generated by tools/make-figures.py -- do not edit by hand.
% Regenerate with: python3 tools/make-figures.py
\begin{tikzpicture}
\begin{axis}[
width=0.60\linewidth, height=4.6cm, tick align=outside, tick pos=left,
  grid=major, grid style={draw=UTPMuted!20},
  axis line style={draw=UTPMuted!70},
  xlabel={Month}, ylabel={Rainfall},
  xmin=0.8, xmax=5.2, ymin=0, ymax=45, xtick={1,2,3,4,5},
  legend style={at={(0.98,0.98)}, anchor=north east, font=\scriptsize},
]
\addplot+[mark=*, mark size=1.8pt, draw=UTPRed, solid] coordinates {(1,7) (2,12) (3,28) (4,3) (5,41)};
\addplot+[mark=*, mark size=1.8pt, draw=UTPBlue, solid] coordinates {(1,14) (2,7) (3,6) (4,19) (5,3)};
\legend{v (red),t (blue)}
\end{axis}
\end{tikzpicture}

\end{utpfig}

\begin{pitfall}
\texttt{lines()} does not expand the axes. Because \texttt{plot(v)} set the
range from \texttt{v} alone, a second series with larger values would be
clipped at the top edge without warning. Set \texttt{ylim} from both vectors
before the first call: \texttt{ylim = c(0, max(v, t))}.
\end{pitfall}

\begin{utpfig}[Choosing a chart]{Start from the question, not from the chart you
find easiest to produce.}{fig:chart-choice}
\begin{tikzpicture}[
  box/.style={draw=DSBlue,fill=DSPaleBlue,rounded corners=2mm,text width=3.0cm,minimum height=1.05cm,align=center,font=\small\bfseries\color{DSNavy}},
  arr/.style={-{Stealth[length=2mm]},thick,draw=DSTeal}
]
\node[box] (question) {What do you want to show?};
\node[box,below left=1.2cm and 2.6cm of question] (part) {Parts of a whole\\Pie};
\node[box,below=1.2cm of question] (dist) {Distribution / groups\\Histogram, box plot};
\node[box,below right=1.2cm and 2.6cm of question] (rel) {Relationship / trend\\Scatter, line};
\draw[arr] (question) -- (part);
\draw[arr] (question) -- (dist);
\draw[arr] (question) -- (rel);
\end{tikzpicture}
\end{utpfig}

Choose a chart that matches the analytical question, then explain the visible evidence.

\begin{quickcheck}
\begin{enumerate}
  \item Which command lists built-in datasets and which command loads \texttt{mtcars}?
  \item How are percentages calculated for the pie chart in the handout?
  \item Which plot compares \texttt{mpg} across \texttt{cyl} groups?
  \item What does \texttt{pairs()} produce in the supplied example?
  \item Which function adds a second series to an existing line plot?
\end{enumerate}
\end{quickcheck}

\begin{chapterrecap}
\begin{itemize}
  \item Match the chart type to composition, distribution, comparison, relationship, or trend.
  \item Use titles, labels, legends, and colors to clarify the evidence.
  \item Every visualization should support a specific analytical claim.
\end{itemize}
\end{chapterrecap}
